% Copyright (c) 2026 Geovana Neves. Licensed under CC BY 4.0 (https://creativecommons.org/licenses/by/4.0/). See LICENSE-AND-AUTHORSHIP.md.
%
\section{One block per rotor}

\begin{frame}[fragile]{One block per rotor: \code{kind = "rotor"}}
\small The block's NAME is an alias over everything that rotor owns.
\begin{lstlisting}[basicstyle=\ttfamily\scriptsize,numbers=none]
[PORT]
kind = "rotor"
x_m = 4.5              # the hub, in the geometry's own frame
y_m = 2.5
z_m = 0.0
axis = "X"              # what it turns about
rpm_sign = 1             # which way
diameter_m = 3.6576      # what an advance ratio resolves against
families_general = []              # hub, spinner: anything not a blade
families_blades = ["Blade1"]       # the blades
blade1 = { azimuth_deg = 0.0, zero = "Y" }   # where blade 1 sits
\end{lstlisting}
\begin{itemize}\small
  \item After this a row says \code{PORT} (\code{MOVING\_BC\_ALIAS} or a
        \code{MOTIONS} record) and states nothing else about that rotor:
        the hub, the axis, the sign, the families and the diameter are all
        here, once.
\end{itemize}
\setupsources{\code{cases/\_\_init\_\_.py}, \code{RotorBlock};
  \code{cases/workflows.py}}
\end{frame}

\begin{frame}{\code{rpm\_sign}: the hand lives in the reference}
\begin{itemize}\small
  \item \code{rpm\_sign} is \code{1} or \code{-1}: which way the rotor
        turns. \code{+1} is the right-hand rule about the block's
        \code{axis}.
  \item A row's own speed, whether stated as \code{RPM} or derived from
        \code{ADVANCE\_RATIO}, is always a \textbf{magnitude}: it says how
        fast, never which way.
  \item A row that restates \code{RPM\_SIGN} beside a rotor the reference
        declares is refused, whether or not the two agree: a row that
        agrees today says nothing when the reference is corrected
        tomorrow, so the refusal is about where the key goes rather than
        about its value.
  \item One exception: a row whose reference declares no rotor block at
        all, using the flat \code{ROTOR\_AXIS} and \code{MOVING\_BOUNDARIES}
        keys, has nowhere else to put the hand, so \code{RPM\_SIGN} beside
        the speed is correct there.
\end{itemize}
\begin{ftslimit}
The hand is declared once, beside the axis, the hub and the blade count it
already declares there; the row and the reference can never disagree
because only one of them is allowed to say it.
\end{ftslimit}
\setupsources{\code{cases/workflows.py}, \code{\_refuse\_a\_row\_restating\_the\_hand},
  \code{\_rpm\_sign}; \code{cases/matrix.py}}
\end{frame}

\begin{frame}{The blade count, and where \code{blade1} sits}
\begin{itemize}\small
  \item \code{families\_general}: what turns with the rotor and is
        \emph{not} a blade, such as the hub or a spinner.
  \item \code{families\_blades}: the blades. \textbf{The blade count is
        the length of this list and nothing else}, so a row states no
        count of its own, and a sector mesh carrying one blade of four
        still reduces over four.
  \item \code{blade1}: \code{azimuth\_deg} (where blade 1 sits, measured
        from \code{zero}) and \code{zero} (the axis that angle is measured
        from). It locates the first family of \code{families\_blades} so
        a per-blade frame can be placed; it does not say which family is
        the first blade, only where that one sits.
  \item A rotor block naming a blade family the opened geometry does not
        carry, or one listing no \code{families\_blades} at all when a row
        needs the count, is refused naming the block and what to add.
\end{itemize}
\setupsources{\code{cases/\_\_init\_\_.py}, \code{RotorBlock};
  \code{cases/workflows.py}}
\end{frame}

\begin{frame}{What else reads the rotor block}
\begin{columns}[T]
\begin{column}{0.50\textwidth}
\begin{block}{The tip and helical Mach numbers}
\small \(R\) is half the block's \code{diameter\_m}, else half the
reference's \code{rotor\_diameter\_m}; a disc's is its
\code{tip\_radius\_m}. With the point's speed \(V\) and speed of sound \(a\):
\(M_\mathrm{tip} = \Omega R / a\) and \(M_\mathrm{hel} =
\sqrt{V^2 + (\Omega R)^2}/a\) [3, 7]. A rotor with no diameter anywhere gets
no Mach numbers, and the plan and the run record say so, naming the row:
nothing is guessed.
\end{block}
\end{column}
\begin{column}{0.46\textwidth}
\begin{block}{A \code{qsteady\_rotor} row}
\small solves the rotor steady, blades still and the free stream turning
about the block's \code{axis} at the row's speed. It is valid only for an
\textbf{isolated, axisymmetric} rotor. The sector case (periodic symmetry)
meshes one blade; the wheel meshes all of them, and
\code{families\_blades} still gives the count \(N\) that spaces the row's
\code{PASSAGE\_POSITIONS} clockings inside one passage, \(360^\circ/N\).
\end{block}
\end{column}
\end{columns}
\setupsources{\code{cases/workflows.py}; the page ``Post-processing
  definitions'', tip and helical Mach numbers}
\end{frame}
